In this section, we first discuss CG applications and the impact of latency on their performance. Then, we review existing unsupervised learning models and window-based approaches used for anomaly detection using time series data.
\subsection{Cloud Gaming applications}
Cloud games are processed and executed on cloud servers and the rendered scenes are streamed over the Internet to client devices removing the need for having powerful gaming computers or consoles \cite{Cai2016ASO}. Several works \cite{marchal2023analysis, Iqbal2021DissectingCG, Domenico2021ANA} focused on studying the behavior of CG platforms and analysed their respective network protocols. For instance, Google Stadia uses WebRTC for its service and relies on RTP to stream its audio and video contents. Marchal et al. \cite{marchal2023analysis} showed that link capacity and latency are driving Stadia platform to adapt its bitrate and resolution.

On the other hand, several works \cite{Claypool2014TheEO,Raaen2014CanGD, Eg2018PlayingWD, Vlahovic2019TheIO} addressed the impact of network latency on the performance of gamers in cloud and non-cloud games. They showed that the performance of gamers drops with the increase of latency. Raaen et al. \cite{Raaen2014CanGD} and Vlahovic et al. \cite{Vlahovic2019TheIO} also include in their subjective studies that most of the gamers could not tolerate a delay above 100ms and the most demanding gamers are able to perceive delays below 40ms. These network perturbations can lead to QoE degradation, which this work aims to detect using unsupervised ML algorithms.

\subsection{Unsupervised Learning models for anomaly detection}
Anomaly detection is a popular and a well-covered research topic with numerous surveys \cite{chandola2009anomaly, Schmidl2022AnomalyDI, Ren2021DeepVA, Cook2020AnomalyDF} proposing different taxonomies and categories of techniques, including existing machine learning models for anomaly detection. Some studies \cite{Pang2020DeepLF, ruff2021unifying, Choi2021DeepLF, 10043819} have specifically focused on the use of deep-learning for anomaly detection due to its efficiency in handling multivariate and high-dimensional data. Schmidl et al. \cite{Schmidl2022AnomalyDI} categorized anomaly detection models for time-series data depending on the way they determine anomalies.
Reconstruction-based algorithms detect anomalies by learning a model from \textit{normal} training data. They encode the training features into a low-dimensional (i.e., latent) space and reconstruct the input features from the latent features. An anomaly score is computed by comparing the reconstructed data to the input data in order to detect anomalies. Since the model is built on \textit{normal} data only, anomalous time-series cannot be well-reconstructed and will have a high anomaly score (i.e., above a pre-determined threshold). This class of algorithms, include Principal Component Analysis (PCA) and neural network algorithms such as AutoEncoder, LSTM-VAE \cite{park2018multimodal}, DAGMM \cite{zong2018deep}, OmniAnomaly \cite{su2019robust}, Donut \cite{xu2018unsupervised}, and USAD \cite{audibert2020usad}. 

One-class classification models detect anomalies by learning a hypersphere that encloses the representation of normal data. Any points that remain outside the learned hypersphere are classified as anomalies. The most popular algorithm from this category is the OC-SVM \cite{scholkopf1999support} and its neural network variant Deep-SVDD \cite{pmlr-v80-ruff18a}.

Furthermore, it is worth mentioning that there are other families of unsupervised learning models for time-series data, including Isolation methods that build an ensemble of isolation trees to isolate anomalies. Isolation methods such as Isolation Forest \cite{liu2008isolation} assume that anomalies are easy to \textit{isolate} since they are fewer in the data and are starkly different from normal instances. Distance-based methods (e.g., KNN \cite{Ramaswamy2000EfficientAF}, LOF\cite{breunig2000lof}) detect anomalous instances based on their larger distances from normal instances. Predictive methods (e.g.,ARIMA) predict time-series sequences and compare them to original time-series to differentiate the anomalies. Generative methods (e.g., GAN \cite{GoodfellowIan2020GenerativeAN}) train a model to generate new (i.e., normal or anomalous) data based on real data distribution and a discriminator that learns how to discriminate real data from generated data. 

The aforementioned anomaly detection models were evaluated in different studies \cite{Nkashama2022RobustnessEO, Garg2022AnEO, Cook2020AnomalyDF} either with broad or domain-specific benchmark datasets. Our work provides a comprehensive and consistent experimental evaluation of anomaly detection models for QoE degradation detection, while taking into account the impact of data contamination and window size.


\subsection{Window-based approaches for anomaly detection}\label{window_metrics}
Chandola et al. \cite{chandola2009anomaly} categorize anomalies into three classes: \textit{point anomalies} which are individual observations that deviate from normal ranges; \textit{collective anomalies} that represent a group of anomalous observations and \textit{contextual anomalies} that represent a group of observations that can be considered as anomalous in a specific context. Point anomalies are the most widely addressed anomalies in the literature. %Citations.

In most time-series, anomalies occur as contiguous anomalous observations (\textit{i.e., collective anomalies}) rather than individual \textit{point anomalies}. Assessing the performance of AD models for time-series data using \ac{PW} approaches is ineffective as it disregards the contiguous nature of anomalies. To overcome this limitation, \ac{PA} approach is proposed by Xu et al. \cite{xu2018unsupervised} and is often employed \cite{su2019robust,audibert2020usad} to deal with windows of anomalies. PA approach considers that all the anomalies of a window are correctly detected by a AD model if any of the anomalous observations in the window is correctly detected. However, it was shown in \cite{Kim2022TowardsAR, wang2023deep, Hwang2022DoYK} that PA approach presents limitations as well and may overestimate the performance of an AD model. The authors showed that a well-trained and efficient algorithm provides the same F1-score as a random model with PA approach. Hence, with the PA, it is difficult to conclude that a model outperforms others. 

To address the limitations of PA approach, \ac{RPA} \cite{hundman2018detecting} and PA$\%K$ \cite{Kim2022TowardsAR} were proposed. Unlike PA, RPA is less tolerant to high false-positive rates by severely penalizing them. Other approaches such as Numenta Anomaly Benchmark (NAB) \cite{Lavin2015EvaluatingRA} and range-based Precision/Recall \cite{Tatbul2018PrecisionAR} have also been proposed, but they are are too complex to be widely adopted. In this chapter, we propose the WAD approach, which addresses the aforementioned limitations and aims at fairly assessing the performance of anomaly detection models.